\begin{musica}{Timor}{João Gil, João Monge}{Tuna Universitária do Porto}
Lavam-se os \ac{E}{o}lhos, ne\ac{A}{ga}-se o \ac{E}{bei}jo,
Do labirinto es\ac{F#7}{co}lhe-se o \ac{B7}{mar};
No cais de\ac{A}{ser}to fica o de\ac{E}{se}jo
Da terra \ac{F#7}{quen}te por \ac{B7}{con}quis\ac{E}{tar}.

Nobre soldado que vens senhor,
Por sobre as asas do teu dragão
Beijas os corpos no chão queimado,
Nunca serás o nosso perdão.

\nota{Refrão:}
\ac{E}{Ai} \ac{B7}{Ti}m\ac{A}{or}
\ac{B7}{Ca}lam-se \ac{E}{as} vozes dos \ac{B7}{teus} a\ac{E}{vós}. \ac{B7}{}__
\ac{C#m}{Ai} \ac{G#}{Ti}m\ac{A}{or}
\ac{B7}{Se} outros \ac{E}{ca}lam can\ac{B7}{te}mos \ac{E}{nós}.

Salgas de ventres que não tiveste,
Ceifando os filhos que não são teus,
Nobre soldado nunca sonhaste
Ver uma espada na mão de Deus.

Da cruz se faz uma lança em chamas
Que sangra o céu no Sol do meio-dia
Do meio dos corpos a mesma lama
Leito final onde o amor nascia.

\nota{Refrão}

\posicoes{C#m: x46654, E: 022100, F#7: 242322, G#: 466544, A: x02220, B7: 021202}
\end{musica}
